\section{System overview \& layering}\label{sec:overview}

QtRocket ships three executables built from one Qt-free simulation core. \code{qtrocket} is
the Qt6 Widgets desktop application; \code{qtrocket-cli} is a headless, line-oriented REPL
over the same simulation API; \code{qtrocket-visualizer} is a standalone OpenGL viewer for
saved designs. All three link the same static library, \code{qtrocket\_core}, which contains
exactly one class --- the \code{QtRocket} application controller --- and pulls in the
\code{model}, \code{sim}, and \code{utils} layer libraries beneath it
(\src{core/CMakeLists.txt}{8}). \Cref{fig:overview-layering} shows the arrangement; the rest of
this section walks it from the top down, then summarizes the build that enforces it.

\tikzfig{system-layering}{The shipped system: three executables over one Qt-free core.
Every edge points strictly downward --- front ends link \texttt{qtrocket\_core},
the core links \texttt{model}/\texttt{sim}/\texttt{utils}, and \texttt{utils} links only
third-party math and networking. The \texttt{Propagatable} interface (right) is the single
seam through which the simulation loop sees the rocket model.}{fig:overview-layering}

\subsection{Three executables, one core}\label{sec:overview-executables}

\begin{description}
\item[\code{qtrocket} (GUI).] The entry point instantiates the \code{Logger} at INFO
  \src{gui/main.cpp}{20}, obtains the controller singleton \src{gui/main.cpp}{24}, then
  constructs \code{QApplication} and \code{MainWindow} directly and blocks in
  \cppinline|a.exec()| \src{gui/main.cpp}{29} --- the conventional Qt main-thread
  arrangement. The executable is the only Qt-dependent front end in the main tree, linking
  \code{Qt6::Widgets}, \code{Qt6::PrintSupport}, and \code{qtrocket\_core}
  \src{gui/CMakeLists.txt}{82}.
\item[\code{qtrocket-cli} (headless REPL).] The CLI entry point deliberately ignores
  \code{argc}/\code{argv} \src{cli/CliMain.cpp}{13}, silences the logger to ERROR because
  stdout doubles as the machine-readable \code{OK}/\code{ERR}/\code{CSV} protocol
  \src{cli/CliMain.cpp}{17}, prints a \code{\#}-prefixed banner, and hands
  \code{std::cin}/\code{std::cout} to \code{cli::Repl} \src{cli/CliMain.cpp}{25}. The Repl
  is compiled into its own \code{cli} static library so that the test binaries drive the
  identical object code users type into \src{cli/CMakeLists.txt}{1}; the executable links that
  library and nothing else \src{cli/CMakeLists.txt}{13}. \Cref{sec:interfaces} reviews the
  command set itself.
\item[\code{qtrocket-visualizer} (OpenGL viewer).] A self-contained Qt application under
  \srcf{visualizer/} that finds its own \code{Qt6::OpenGLWidgets}/\code{OpenGL} components
  \src{visualizer/CMakeLists.txt}{11} and, by stated design, never touches \srcf{gui/} or
  \srcf{cli/} code \src{visualizer/CMakeLists.txt}{3}. Its \code{main} disables the RHI
  widget backing store \src{visualizer/main.cpp}{22} and deliberately requests no specific
  GL version (a documented EGL/Wayland-plus-NVIDIA workaround,
  \src{visualizer/main.cpp}{10}), then loads an optional \code{.qrd} design path from
  \code{argv[1]} \src{visualizer/main.cpp}{35} through the shared
  \code{model::DesignSerializer} \src{visualizer/VisualizerWindow.cpp}{184}.
\end{description}

The visualizer's load path illustrates the document's recurring motif early: the shared
design serializer was consumed by the CLI and the visualizer for months before the GUI ---
the front end a desktop user would expect to open files from --- gained its File\,>\,Open/%
Save wiring with the \code{RocketTreeViewImplementation} merge
(\cref{sec:interfaces,sec:persistence}). The seam was built and tested well before its most
obvious consumer arrived.

\subsection{The layering and the link graph}\label{sec:overview-layers}

The layering is not a documentation convention; it is the link graph. \code{utils} sits at
the bottom, linking only \code{libcurl} and \code{eigen}
\src{core/utils/CMakeLists.txt}{17}. \code{sim} links only \code{utils}
\src{core/sim/CMakeLists.txt}{31}. \code{model} links \code{utils}, \code{sim},
\code{Boost::property\_tree}, and \code{jsoncpp\_static} \src{core/model/CMakeLists.txt}{47} ---
so \code{model} and \code{sim} are layered rather than peers: the model depends on the
simulation layer's value types (\code{sim::Aero}, \code{StateData}), never the reverse
\src{core/model/CMakeLists.txt}{44}. Above both sits \code{qtrocket\_core}
(\cref{lst:overview-core}), and above that the three executables plus the test binaries.
Include resolution everywhere uses a single global
\cppinline|include_directories(${PROJECT_SOURCE_DIR})| \src{CMakeLists.txt}{122}, which is
what makes the \code{model/...}, \code{sim/...}, and \code{utils/...} path prefixes resolve
uniformly across the tree.

\begin{listing}[htbp]
\begin{minted}{text}
# Application controller shared by the GUI, the CLI, and the
# integration tests.
add_library(qtrocket_core STATIC
   QtRocket.cpp
   QtRocket.h)
qtrocket_enable_coverage(qtrocket_core)

target_link_libraries(qtrocket_core PUBLIC
   model
   sim
   utils)
\end{minted}
\caption{The keystone of the link graph \srccap{core/CMakeLists.txt}{6}: the controller
compiles once into a static library, and every executable and integration-test binary links
it rather than recompiling \texttt{QtRocket.cpp}.}\label{lst:overview-core}
\end{listing}

\begin{designrule}[The layering invariant, stated where it binds]
``Bottom layer: utils must not depend on model/ or sim/'' \src{core/utils/CMakeLists.txt}{16}.
The comment sits directly above the \cppinline|target_link_libraries| call it constrains,
and the review confirmed it holds: no header under \srcf{utils/} includes anything from
\srcf{model/} or \srcf{sim/}. This single rule is what keeps the dependency graph acyclic
--- every layer above utils may share its math typedefs and logging without creating a
back-edge.
\end{designrule}

\subsection{The controller and the \code{Propagatable} bridge}\label{sec:overview-controller}

\code{QtRocket} is the application controller: a process-wide singleton
\src{core/QtRocket.cpp}{19} that owns the \code{(RocketModel, Propagator)} pair, the
\code{Environment}, and the \code{MotorModelDatabase}. Installing a rocket constructs its
propagator in the same breath, with the shared environment injected
\src{core/QtRocket.h}{41} --- the model and its integrator loop are never allowed to drift apart.
The API surface the front ends drive is deliberately small: accessors for the environment,
time step, integrator selection, rocket, and motor database \src{core/QtRocket.h}{34};
\code{launchRocket()} \src{core/QtRocket.h}{43}; the recorded state history via
\code{getStates()} \src{core/QtRocket.h}{48}; whole-flight summary statistics
\src{core/QtRocket.h}{55}; a typed termination reason \src{core/QtRocket.h}{62}; and
\code{setInitialState()} \src{core/QtRocket.h}{68}. A launch is four steps --- clear the state
history, reset propagator time to zero, ignite the motor, run the ODE loop to termination
\src{core/QtRocket.cpp}{54} --- and everything else (staging launch parameters, choosing models,
exporting results) lives in the front ends.

\begin{keyinsight}[Racy double-checked initialization]
\code{QtRocket::getInstance()} reads the static \code{initialized} flag without holding the
mutex \src{core/QtRocket.cpp}{21}; only \code{init()} re-checks under the lock
\src{core/QtRocket.cpp}{30}. Two threads calling \code{getInstance()} concurrently would race on
that unsynchronized read --- a data race in the C++ memory model. It is latent today because
the two front ends that use it call it once, single-threaded, at startup (the visualizer
builds its own \code{RocketModel} and never touches the controller); it becomes real the
moment simulation moves to the worker thread that \cref{sec:roadmap} anticipates. The
\code{Logger} singleton one layer down already uses the safe Meyers form
\src{utils/Logger.cpp}{19}.
\end{keyinsight}

The seam between the two halves of the core is the \code{Propagatable} interface
\src{model/Propagatable.h}{23}: the only view of the rocket that the simulation loop ever
sees. \code{RocketModel} implements it; the \code{Propagator} consumes it. Its pure-virtual
heart is four questions the integrator asks at every step
(\cref{lst:overview-propagatable}), plus \code{terminateCondition()}
\src{model/Propagatable.h}{39}, which lets the model rather than the loop decide when a
flight is over. Only the three linear degrees of freedom are integrated today, so the force
half of the bridge does all the work; what the torque half is waiting for is reviewed with
the aerodynamics in \cref{sec:aero}, and the loop that drives the interface in
\cref{sec:sim}. The mass-property machinery behind \code{getMass(t)} --- the composite part
tree --- is the subject of \cref{sec:model}.

\begin{listing}[htbp]
\begin{minted}{cpp}
   virtual Vector3 getForces(double t, const Vector3& position, const Vector3& velocity, sim::Environment& environment) = 0;
   virtual Vector3 getTorques(double t) = 0;

   virtual double getMass(double t) = 0;
   virtual Matrix3 getCompositeInertiaTensor(double t) = 0;
\end{minted}
\caption{The model$\leftrightarrow$sim bridge \srccap{model/Propagatable.h}{29}: the
integrator needs forces, torques, and time-varying mass properties, and nothing else about
the rocket.}\label{lst:overview-propagatable}
\end{listing}

\subsection{Build system in brief}\label{sec:overview-build}

The build is a single CMake project --- \cppinline|cmake_minimum_required(VERSION 3.25)|
\src{CMakeLists.txt}{1}, \code{VERSION 0.1} \src{CMakeLists.txt}{3} --- pinned to C++23
\src{CMakeLists.txt}{5} with \code{compile\_commands.json} exported for tooling
\src{CMakeLists.txt}{7}. Generator, compiler, and build type are pinned by
\srcf{CMakePresets.json}: Ninja everywhere \src{CMakePresets.json}{7}, with eight configure
presets --- \code{debug}/\code{release} for the default compiler, Clang, and GCC, a
release-only MSVC preset \src{CMakePresets.json}{70}, and \code{coverage-clang}, which
redirects to a separate \code{build-coverage/} tree \src{CMakePresets.json}{35} so
instrumentation never disturbs the normal build. Matching build and test presets exist for
each, and the same presets drive both developer workflows and CI.

\begin{table}[htbp]
\centering\small
\begin{tabularx}{\linewidth}{@{}ll>{\raggedright\arraybackslash}Xl@{}}
\toprule
Dependency & Pinned tag & Build shape and role & \\
\midrule
GoogleTest & \code{v1.17.0} & test framework for all six test binaries & \src{CMakeLists.txt}{16} \\
jsoncpp & \code{1.9.7} & static-only, tests/examples off; thrustcurve.org JSON & \src{CMakeLists.txt}{31} \\
curl & \code{curl-8\_20\_0} & static, HTTP-only, optional backends forced off & \src{CMakeLists.txt}{43} \\
Eigen & \code{5.0.1} & vectors, matrices, quaternions beneath \code{utils} & \src{CMakeLists.txt}{66} \\
Boost & \code{boost-1.91.0-1} & restricted to \code{property\_tree}; XML persistence & \src{CMakeLists.txt}{73} \\
Qt6 & system & Widgets/PrintSupport; visualizer adds OpenGL modules & \src{gui/CMakeLists.txt}{13} \\
\bottomrule
\end{tabularx}
\caption{Every non-Qt dependency is vendored via FetchContent at an exact tag; Qt6 is the
sole system dependency. The first configure therefore downloads and compiles the world ---
slow once, hermetic thereafter.}\label{tab:overview-deps}
\end{table}

Dependency policy is strict vendoring (\cref{tab:overview-deps}). The curl configuration is
the most deliberate: built static and HTTP-only \src{CMakeLists.txt}{46} with every optional
third-party backend explicitly forced off, because auto-detected Homebrew libraries on macOS
CI runners had previously broken the final link --- the hermeticity rationale is recorded in
the tree \src{CMakeLists.txt}{48}. Qt is resolved with a hard
\cppinline|find_package(Qt6 REQUIRED ...)| \src{gui/CMakeLists.txt}{13}; the finer
hygiene findings in that block are examined in \cref{sec:testing}.

Warning hygiene is tree-wide and strict: \code{-Wall -Wextra -Wpedantic -Werror} on GCC and
Clang \src{CMakeLists.txt}{91}, \code{/W4 /WX} on MSVC \src{CMakeLists.txt}{82}, applied
after the FetchContent blocks so vendored code is exempt. The test suites themselves (six
GoogleTest binaries), the five-platform CI matrix that runs them, the llvm-cov coverage
target, and the finer warning-hygiene findings (a stale sanitizer claim, the vendored
\srcf{gui/qcustomplot.cpp} exemption) are reviewed in depth in \cref{sec:testing}.

Finally, what the build does \emph{not} yet do: distribution. QtRocket today builds well
everywhere and installs properly nowhere --- the install rule, bundle metadata, and
translation rules are template stubs, examined with the rest of the packaging story in
\cref{sec:testing}; \cref{sec:roadmap} places that gap among the others.
